\begingroup
\setlength{\leftskip}{0pt}
\setlength{\rightskip}{0pt}
\setlength{\parfillskip}{0pt plus 1fil}
\setlength{\parindent}{0pt}
\setlength{\parskip}{2pt}
\movement{TECHNICAL NOTES}{Looking more closely}

These notes examine the nine-pixel experiment's geometry, derivatives and reproducibility, then give the ring experiment's mathematical and reproduction details. Both examples use synthetic data generated by code, not measurements from a physical experiment. Their training runs and reported results are computed, while the explicitly labelled hand-built constructions are not~learned.

\input{chapters/v47-01-neural-networks/reader-lab-v47}


\Needspace{65mm}
\smallhead{Why one affine score is insufficient}

Let $r_1,r_2,r_3$ be the three horizontal templates and $q_1,q_2,q_3$ the vertical ones, represented as nine-dimensional vectors. Then
\[
\sum_{k=1}^3 r_k=\sum_{k=1}^3 q_k=\boldsymbol{1}.
\]
Here $\boldsymbol{1}$ is the nine-entry vector with every entry equal to one. Each pixel belongs to exactly one horizontal template and exactly one vertical template. For any affine score $s(x)=u^\top x+d$, it follows that
\[
\sum_{k=1}^3s(r_k)=u^\top\boldsymbol{1}+3d
=\sum_{k=1}^3s(q_k).
\]
Our decision convention requires $s(r_k)<0$ for all horizontal templates and $s(q_k)\geq0$ for all vertical templates. The first sum would be negative and the second nonnegative, contradicting their equality. A different scalar threshold can be absorbed into $d$.

Six nonlinear hidden units are sufficient, although this does not prove six are necessary. For a constructive example, assign one unit to each row and column. Give its three relevant pixels weight 1, the others weight 0, and use bias $-2.5$. A complete matching stroke activates its detector to 0.5. A perpendicular stroke contributes only 1 to each such sum and activates none of those detectors. Subtract the sum of row-detector outputs from the sum of column-detector outputs. The clean horizontal templates score $-0.5$ and the vertical templates $+0.5$.

Those are hand-constructed parameters proving capacity on the six clean templates. They are \emph{not} the parameters learned in the reported run, and the construction does not establish noise~performance.



\smallhead{Where the compact gradient comes from}
\label{note:v29-gradient}

First differentiate the loss with respect to the probability, then the probability with respect to the score:
\[
\begin{aligned}
\frac{\partial\ell}{\partial p}
 &=-\frac{y}{p}+\frac{1-y}{1-p}
 =\frac{p-y}{p(1-p)},\\
\frac{\mathrm dp}{\mathrm dz}&=p(1-p).
\end{aligned}
\]
For finite $z$, sigmoid has $0<p<1$, so these denominators are nonzero. Multiply the two derivatives. Their $p(1-p)$ factors cancel, leaving
\[
\frac{\partial\ell}{\partial z}=p-y.
\]
The chain from $w_{ij}$ to the loss then passes through $a_j$, $h_j$ and $z$:
\[
\begin{aligned}
\frac{\partial\ell}{\partial w_{ij}}&=
\frac{\partial\ell}{\partial z}
\frac{\partial z}{\partial h_j}
\frac{\partial h_j}{\partial a_j}
\frac{\partial a_j}{\partial w_{ij}}
\\&=(p-y)v_j\mathbf{1}[a_j>0]x_i.
\end{aligned}
\]
\Needspace{4\baselineskip}
This is an ordinary derivative when $a_j\ne0$. At the ReLU kink, $a_j=0$, the program uses zero for that gate's derivative, a computational convention rather than an ordinary derivative of~ReLU.
Average these derivatives over the training batch to obtain the derivative of mean loss. Differentiation is linear under this finite average. Backpropagation efficiently reuses intermediate derivatives rather than perturbing every parameter~separately.

For stable loss calculation, the experiment uses the equivalent expression
\[
\ell(z,y)=\log(1+e^z)-yz,
\]
evaluating its first term with \texttt{logaddexp(0, z)}. This avoids forming the logarithm of a probability that has rounded to zero or~one.

\smallhead{Input sensitivity and parameter symmetry}
\label{note:v31-positive-rescaling}

Away from ReLU gate boundaries, the score's sensitivity to input $i$ is
\[
\frac{\partial z}{\partial x_i}
=\sum_j v_j\mathbf{1}[a_j>0]w_{ij}.
\]
Crossing an activation boundary changes this formula's active terms. Thus a derivative at one picture is not automatically an explanation of a large change to~it.

There is also a parameter symmetry. Multiply all incoming weights and the bias of a hidden unit by a positive number $\lambda$, and divide its outgoing weight by $\lambda$. Because $\max(0,\lambda a)=\lambda\max(0,a)$, the network's output is unchanged for every input. Hidden units can also be permuted with matching permutations of their connections. Different parameter descriptions can therefore implement exactly the same~function.

This is one concrete reason to resist treating the magnitude of a single weight as an intrinsic measure of meaning. For our purposes the relevant object is often the computation and its response to specified interventions, not a preferred numbering or scale of internal~units. Equal forward functions need not remain equal after ordinary gradient descent: Chapter~\ref{ch:v8-2}'s notes give an exact one-update counterexample on page~\pageref{note:v29-equivalent-updates}.



\Needspace{13\baselineskip}
\smallhead{Reproducing and delimiting the experiment}

The recorded run used NumPy 2.3.5, with seeds 11 for initialization, 17 for training data and 29 for test data. Initial input and output weights are zero-mean normal draws with standard deviation 0.4. Independent pixel noise has standard deviation 0.12 before clipping into $[0,1]$. Hidden biases start at 0.1 and the output bias at zero. Updates use the full batch, not sampled mini-batches. The historical producer is \path{examples/ch01/strokes_v3.py}. For read-only exploration use the ZIP's \texttt{investigation\_v9.py} route above.

The output folder stores the datasets, initial and final parameters, training curve, predictions, numerical checks and configuration. Evaluation occurs after the fixed training schedule. There was no hyperparameter search against the test results. This is one initialization and one data split. It does not measure variability across training runs. The derivative check perturbs each initial parameter by $10^{-6}$ in both directions and compares the central finite difference with backpropagation.

The 600 test pictures come from six repeated template families. Treating them as 600 unrelated kinds of visual challenge would overstate the evidence. The contrast and deletion evaluations are small, deterministic checks, not estimates of a general out-of-distribution success rate. All results in this chapter are illustrative computations, not new research claims about neural-network performance.

\input{chapters/v47-01-neural-networks/ring-technical-v47}


\Needspace{6\baselineskip}
\smallhead{The tools behind the diagrams}

The NumPy program exposes the operations and derivatives. These four frameworks express the same network through different notation. This is a short tour, not a software survey.

\noindent\begin{minipage}{\linewidth}
For our network, the central forward calculation is only four lines after the NumPy import. This compact version shows the algebra:
\begin{lstlisting}[language=Python]
import numpy as np
a = x @ W + b
h = np.maximum(a, 0)
z = h @ v + c
p = 1 / (1 + np.exp(-z))
\end{lstlisting}
\end{minipage}
The symbol \texttt{@} means matrix multiplication. For a batch of $N$ pictures, \texttt{x} has shape $N\times9$, \texttt{W} has shape $9\times6$, and the hidden values have shape $N\times6$. The bias $b$ and output weights $v$ each have six entries, and $c$ is a scalar. The last weighted sum produces one score per~picture. The runnable reader program above uses a numerically stable sigmoid for extreme scores.

Each version implements Figure~\ref{fig:forward} up to score $z$. Reproducing its displayed probability also requires sigmoid and the same trained weights and~biases.

An architecture specifies how the numbers are combined. A trained parameter set specifies which particular calculation that architecture performs. Library defaults can give the same architecture a different starting~point.

\noindent\begin{minipage}{\linewidth}
PyTorch expresses the same architecture as:
\begin{lstlisting}[language=Python]
import torch

model = torch.nn.Sequential(
    torch.nn.Linear(9, 6),
    torch.nn.ReLU(),
    torch.nn.Linear(6, 1),
)
\end{lstlisting}
\end{minipage}

\noindent\begin{minipage}{\linewidth}
TensorFlow spells out the weighted sums, with parameters supplied explicitly:
\begin{lstlisting}[language=Python]
import tensorflow as tf

def tensorflow_score(x, params):
    W, b, V, c = params
    h = tf.nn.relu(tf.linalg.matmul(x, W) + b)
    return tf.linalg.matmul(h, V) + c
\end{lstlisting}

Here \texttt{params} contains four floating-point tensors: \texttt{W} has shape $9\times6$, \texttt{b} has shape $(6,)$, \texttt{V} has shape $6\times1$, and \texttt{c} has shape $(1,)$. The input has shape $N\times9$. The result has shape $N\times1$, one raw score per picture. All inputs and parameters must use compatible data types. The capital \texttt{V} holds the earlier six-element output weight vector as a~column.
\end{minipage}

\noindent\begin{minipage}{\linewidth}
JAX expresses the same calculation as a function of parameters and input:
\begin{lstlisting}[language=Python]
import jax
import jax.numpy as jnp

def jax_score(params, x):
    W, b, V, c = params
    h = jax.nn.relu(jnp.matmul(x, W) + b)
    return jnp.matmul(h, V) + c
\end{lstlisting}
\end{minipage}

Supply JAX arrays with the same shapes. This function returns scores without changing the supplied~parameters.

\noindent\begin{minipage}{\linewidth}
Keras 3 provides a higher-level layer API over these backends. Here we select TensorFlow before importing Keras:
\begin{lstlisting}[language=Python]
import os
os.environ["KERAS_BACKEND"] = "tensorflow"
import keras

model = keras.Sequential([
    keras.Input(shape=(9,)),
    keras.layers.Dense(6, activation="relu"),
    keras.layers.Dense(1),
])
\end{lstlisting}
\end{minipage}

Run this example in a fresh Python session with Keras 3 and the chosen backend installed. To use JAX or PyTorch instead, set the value to \texttt{"jax"} or \texttt{"torch"} before the first Keras import. The layer definition stays the same. Keras creates its parameters, as the PyTorch layers~do.

For training, a logits-aware binary cross-entropy loss combines sigmoid and loss calculation stably. Do not apply sigmoid first when the loss expects logits. These are forward examples, not complete training programs. The TensorFlow and JAX functions omit parameter initialization.

Reproducing training requires matching the data, initial parameters, loss reduction and update rule. PyTorch's weight matrices transpose our NumPy convention. All four forward examples were checked against NumPy on 600 saved test images with copied parameters. Matching classifications does not mean identical scores: floating-point precision and intermediate conversions matter. The companion records backend versions, observed data types and numerical tolerances. These are forward checks, not training replications. The reported training measurements come from NumPy.

The frameworks handle the layer calculations and derivatives. We still choose the labels, the loss and the evaluation~data.





\clearpage
\fontsize{10}{12.5}\selectfont
\smallhead{Sources and further reading}

The companion implementation supplies the equations and numerical traces. The official documentation below supports the framework operations and mappings. Documentation checked September 2026.
\begin{enumerate}
\item D. E. Rumelhart, G. E. Hinton and R. J. Williams, \href{https://doi.org/10.1038/323533a0}{Learning representations by back-propagating errors}, \emph{Nature} 323 (1986), 533--536. A primary account of learning hidden representations with derivatives propagated through a network.
\item C. Guo, G. Pleiss, Y. Sun and K. Q. Weinberger, \href{https://proceedings.mlr.press/v70/guo17a.html}{On Calibration of Modern Neural Networks}, PMLR 70 (2017), 1321--1330. This supplies background for the distinction between accurate classifications and reliable probabilities. It does not evaluate our small experiment.
\item J. Hawkins and S. Ahmad, \href{https://doi.org/10.3389/fncir.2016.00023}{Why Neurons Have Thousands of Synapses, a Theory of Sequence Memory in Neocortex}, \emph{Frontiers in Neural Circuits} 10 (2016), 23. A proposed dendritic sequence-memory model, not evidence that ordinary backpropagation is how every biological neuron learns. Chapter~\ref{ch:v8-3} returns to the separate issue of retaining context across observations.
\item PyTorch, \href{https://docs.pytorch.org/docs/2.14/generated/torch.nn.Linear.html}{\texttt{Linear}}: affine transformation and weight layout.
\item PyTorch, \href{https://docs.pytorch.org/docs/2.14/generated/torch.nn.BCEWithLogitsLoss.html}{\texttt{BCEWithLogitsLoss}}: a stable combination of sigmoid and binary cross-entropy.
\item PyTorch, \href{https://docs.pytorch.org/tutorials/beginner/basics/autogradqs_tutorial.html}{Automatic differentiation}: computation graphs and gradient calculation.
\item Keras, \href{https://keras.io/api/layers/core_layers/dense/}{\texttt{Dense}}: weighted sums, biases and activation functions.
\item Keras, \href{https://keras.io/keras_3/}{Keras 3}: the relationship between its high-level API and its backends.
\item Keras, \href{https://keras.io/getting_started/}{Getting started}: selecting the backend before importing Keras.
\item TensorFlow, \href{https://www.tensorflow.org/api_docs/python/tf/linalg/matmul}{\texttt{tf.linalg.matmul}} and \href{https://www.tensorflow.org/api_docs/python/tf/nn/relu}{\texttt{tf.nn.relu}}: matrix products and rectification.
\item JAX, \href{https://docs.jax.dev/en/latest/_autosummary/jax.numpy.matmul.html}{\texttt{jax.numpy.matmul}} and \href{https://docs.jax.dev/en/latest/_autosummary/jax.nn.relu.html}{\texttt{jax.nn.relu}}: the corresponding array operations.
\end{enumerate}
\par\endgroup
